\documentclass[12pt,a4paper]{article}
\usepackage[UTF8]{ctex}
\usepackage{amsmath,amssymb,amsthm}
\usepackage{geometry}

\geometry{margin=2.5cm}

\title{零空间和行空间的详细解释}
\author{第11章补充说明}
\date{}

\begin{document}

\maketitle

\section{问题提出}

在解释约束力$F_{\text{constraint}} = A^T(\theta)\lambda$时，我们提到了：
\begin{itemize}
\item 约束子空间：由$A(\theta)$的\textbf{零空间}定义
\item 约束力子空间：由$A^T(\theta)$的\textbf{列空间}定义（即$A(\theta)$的\textbf{行空间}）
\end{itemize}

\textbf{问题}：
\begin{enumerate}
\item 什么是零空间（null space）？
\item 什么是行空间（row space）？
\item 它们之间有什么关系？
\item 在约束力问题中如何应用？
\end{enumerate}

\section{零空间（Null Space）}

\subsection{定义}

\textbf{零空间}：对于矩阵$A \in \mathbb{R}^{m \times n}$，其零空间定义为：
\begin{equation}
\text{Null}(A) = \{x \in \mathbb{R}^n : Ax = 0\}
\end{equation}

\textbf{通俗理解}：
\begin{itemize}
\item 零空间是所有使得$Ax = 0$的向量$x$的集合
\item 这些向量被矩阵$A$"映射"到零向量
\item 零空间是$\mathbb{R}^n$的一个子空间
\end{itemize}

\subsection{性质}

\textbf{性质1：零空间是子空间}
\begin{itemize}
\item 如果$x_1, x_2 \in \text{Null}(A)$，那么$x_1 + x_2 \in \text{Null}(A)$
\item 如果$x \in \text{Null}(A)$，$c$是标量，那么$cx \in \text{Null}(A)$
\end{itemize}

\textbf{性质2：零空间的维度}
\begin{itemize}
\item 如果$A$的秩为$r$，那么$\text{Null}(A)$的维度为$n - r$
\item 这称为"秩-零度定理"（Rank-Nullity Theorem）
\end{itemize}

\textbf{性质3：零空间与行空间正交}
\begin{itemize}
\item 零空间$\text{Null}(A)$与行空间$\text{Row}(A)$是正交的
\item 即：如果$x \in \text{Null}(A)$，$y \in \text{Row}(A)$，那么$x^T y = 0$
\end{itemize}

\subsection{具体例子}

\textbf{例子1：简单矩阵}

设$A = \begin{bmatrix} 1 & 0 & 0 \\ 0 & 1 & 0 \end{bmatrix} \in \mathbb{R}^{2 \times 3}$

\textbf{零空间}：
\begin{align}
Ax &= 0 \\
\begin{bmatrix} 1 & 0 & 0 \\ 0 & 1 & 0 \end{bmatrix} \begin{bmatrix} x_1 \\ x_2 \\ x_3 \end{bmatrix} &= \begin{bmatrix} 0 \\ 0 \end{bmatrix} \\
x_1 &= 0 \\
x_2 &= 0
\end{align}

因此：
\begin{equation}
\text{Null}(A) = \left\{\begin{bmatrix} 0 \\ 0 \\ x_3 \end{bmatrix} : x_3 \in \mathbb{R}\right\} = \text{span}\left\{\begin{bmatrix} 0 \\ 0 \\ 1 \end{bmatrix}\right\}
\end{equation}

\textbf{维度}：
\begin{itemize}
\item $A$的秩：$r = 2$
\item 零空间的维度：$n - r = 3 - 2 = 1$
\item 零空间由向量$\begin{bmatrix} 0 \\ 0 \\ 1 \end{bmatrix}$张成
\end{itemize}

\textbf{例子2：约束矩阵}

设$A(\theta) = \begin{bmatrix} 1 & 0 & 0 & 0 & 0 & 0 \\ 0 & 1 & 0 & 0 & 0 & 0 \\ 0 & 0 & 0 & 0 & 0 & 1 \end{bmatrix} \in \mathbb{R}^{3 \times 6}$

\textbf{约束}：$A(\theta)V = 0$，其中$V = \begin{bmatrix} \omega_x \\ \omega_y \\ \omega_z \\ v_x \\ v_y \\ v_z \end{bmatrix}$

\textbf{零空间}：
\begin{align}
A(\theta)V &= 0 \\
\begin{bmatrix} \omega_x \\ \omega_y \\ v_z \end{bmatrix} &= \begin{bmatrix} 0 \\ 0 \\ 0 \end{bmatrix}
\end{align}

因此：
\begin{equation}
\text{Null}(A(\theta)) = \left\{\begin{bmatrix} 0 \\ 0 \\ \omega_z \\ v_x \\ v_y \\ 0 \end{bmatrix} : \omega_z, v_x, v_y \in \mathbb{R}\right\}
\end{equation}

\textbf{物理意义}：
\begin{itemize}
\item 零空间表示所有满足约束的速度$V$
\item 即：$\omega_x = 0, \omega_y = 0, v_z = 0$
\item 可以自由运动的方向：$\omega_z, v_x, v_y$（3个自由度）
\end{itemize}

\section{行空间（Row Space）}

\subsection{定义}

\textbf{行空间}：对于矩阵$A \in \mathbb{R}^{m \times n}$，其行空间定义为：
\begin{equation}
\text{Row}(A) = \{A^T y : y \in \mathbb{R}^m\} = \text{Col}(A^T)
\end{equation}

\textbf{等价定义}：
\begin{itemize}
\item 行空间是$A$的所有行的线性组合的集合
\item 行空间是$A^T$的列空间
\item 行空间是$\mathbb{R}^n$的一个子空间
\end{itemize}

\subsection{性质}

\textbf{性质1：行空间是子空间}
\begin{itemize}
\item 行空间是$\mathbb{R}^n$的一个子空间
\item 维度等于$A$的秩$r$
\end{itemize}

\textbf{性质2：行空间与零空间正交}
\begin{itemize}
\item 行空间$\text{Row}(A)$与零空间$\text{Null}(A)$是正交的
\item 即：如果$x \in \text{Row}(A)$，$y \in \text{Null}(A)$，那么$x^T y = 0$
\end{itemize}

\textbf{性质3：行空间与列空间的关系}
\begin{itemize}
\item 行空间的维度 = 列空间的维度 = $A$的秩$r$
\item 行空间$\text{Row}(A)$ = 列空间$\text{Col}(A^T)$
\end{itemize}

\subsection{具体例子}

\textbf{例子1：简单矩阵}

设$A = \begin{bmatrix} 1 & 0 & 0 \\ 0 & 1 & 0 \end{bmatrix} \in \mathbb{R}^{2 \times 3}$

\textbf{行空间}：
\begin{itemize}
\item 行1：$\begin{bmatrix} 1 & 0 & 0 \end{bmatrix}$
\item 行2：$\begin{bmatrix} 0 & 1 & 0 \end{bmatrix}$
\end{itemize}

行空间是所有行的线性组合：
\begin{equation}
\text{Row}(A) = \left\{c_1\begin{bmatrix} 1 \\ 0 \\ 0 \end{bmatrix} + c_2\begin{bmatrix} 0 \\ 1 \\ 0 \end{bmatrix} : c_1, c_2 \in \mathbb{R}\right\} = \text{span}\left\{\begin{bmatrix} 1 \\ 0 \\ 0 \end{bmatrix}, \begin{bmatrix} 0 \\ 1 \\ 0 \end{bmatrix}\right\}
\end{equation}

\textbf{维度}：
\begin{itemize}
\item $A$的秩：$r = 2$
\item 行空间的维度：$r = 2$
\end{itemize}

\textbf{例子2：约束矩阵}

设$A(\theta) = \begin{bmatrix} 1 & 0 & 0 & 0 & 0 & 0 \\ 0 & 1 & 0 & 0 & 0 & 0 \\ 0 & 0 & 0 & 0 & 0 & 1 \end{bmatrix} \in \mathbb{R}^{3 \times 6}$

\textbf{行空间}：
\begin{itemize}
\item 行1：$\begin{bmatrix} 1 & 0 & 0 & 0 & 0 & 0 \end{bmatrix}$（对应$\omega_x$方向）
\item 行2：$\begin{bmatrix} 0 & 1 & 0 & 0 & 0 & 0 \end{bmatrix}$（对应$\omega_y$方向）
\item 行3：$\begin{bmatrix} 0 & 0 & 0 & 0 & 0 & 1 \end{bmatrix}$（对应$v_z$方向）
\end{itemize}

行空间：
\begin{equation}
\text{Row}(A(\theta)) = \text{span}\left\{\begin{bmatrix} 1 \\ 0 \\ 0 \\ 0 \\ 0 \\ 0 \end{bmatrix}, \begin{bmatrix} 0 \\ 1 \\ 0 \\ 0 \\ 0 \\ 0 \end{bmatrix}, \begin{bmatrix} 0 \\ 0 \\ 0 \\ 0 \\ 0 \\ 1 \end{bmatrix}\right\}
\end{equation}

\textbf{物理意义}：
\begin{itemize}
\item 行空间表示约束方向
\item 即：$\omega_x$方向、$\omega_y$方向、$v_z$方向
\item 约束力必须位于这些方向上
\end{itemize}

\section{零空间与行空间的关系}

\subsection{正交性}

\textbf{定理}：零空间$\text{Null}(A)$与行空间$\text{Row}(A)$是正交的。

\textbf{证明}：

设$x \in \text{Null}(A)$，$y \in \text{Row}(A)$。

由于$y \in \text{Row}(A)$，存在$z \in \mathbb{R}^m$使得$y = A^T z$。

由于$x \in \text{Null}(A)$，有$Ax = 0$。

计算内积：
\begin{align}
x^T y &= x^T (A^T z) \\
&= (Ax)^T z \\
&= 0^T z = 0
\end{align}

因此，$x^T y = 0$，即零空间与行空间正交。$\square$

\subsection{互补性}

\textbf{秩-零度定理}：
\begin{equation}
\dim(\text{Null}(A)) + \dim(\text{Row}(A)) = n
\end{equation}

其中$n$是$A$的列数。

\textbf{几何意义}：
\begin{itemize}
\item 零空间和行空间将$\mathbb{R}^n$分解为两个正交的子空间
\item 零空间的维度 + 行空间的维度 = $n$
\item 任何向量$x \in \mathbb{R}^n$可以唯一分解为：
\begin{equation}
x = x_{\text{null}} + x_{\text{row}}
\end{equation}
其中$x_{\text{null}} \in \text{Null}(A)$，$x_{\text{row}} \in \text{Row}(A)$
\end{itemize}

\subsection{具体例子}

\textbf{例子：约束矩阵$A(\theta) \in \mathbb{R}^{3 \times 6}$}

\textbf{维度}：
\begin{itemize}
\item $A$的秩：$r = 3$
\item 零空间的维度：$6 - 3 = 3$
\item 行空间的维度：$3$
\item 验证：$3 + 3 = 6$ ✓
\end{itemize}

\textbf{分解}：
\begin{itemize}
\item 任何6维向量$V$可以分解为：
\begin{equation}
V = V_{\text{null}} + V_{\text{row}}
\end{equation}
\item $V_{\text{null}} \in \text{Null}(A)$：满足约束的速度（自由运动方向）
\item $V_{\text{row}} \in \text{Row}(A)$：违反约束的速度（约束方向）
\end{itemize}

\section{在约束力问题中的应用}

\subsection{约束子空间}

\textbf{定义}：约束子空间是速度$V$必须位于的子空间。

\textbf{数学表述}：
\begin{equation}
\text{约束子空间} = \text{Null}(A(\theta))
\end{equation}

\textbf{物理意义}：
\begin{itemize}
\item 所有满足约束$A(\theta)V = 0$的速度$V$
\item 这些速度可以自由运动，不受约束力影响
\item 维度：$6 - k$（$k$是约束数量）
\end{itemize}

\subsection{约束力子空间}

\textbf{定义}：约束力$F_{\text{constraint}}$必须位于的子空间。

\textbf{数学表述}：
\begin{equation}
\text{约束力子空间} = \text{Row}(A(\theta)) = \text{Col}(A^T(\theta))
\end{equation}

\textbf{物理意义}：
\begin{itemize}
\item 约束力必须位于约束方向
\item 约束力不会影响自由运动方向
\item 维度：$k$（$k$是约束数量）
\end{itemize}

\subsection{约束力的形式}

\textbf{约束力}：
\begin{equation}
F_{\text{constraint}} = A^T(\theta)\lambda
\end{equation}

\textbf{为什么？}
\begin{itemize}
\item $A^T(\theta)$的列空间 = $\text{Row}(A(\theta))$
\item 因此$F_{\text{constraint}} = A^T(\theta)\lambda \in \text{Row}(A(\theta))$
\item 这确保了约束力位于约束力子空间中
\end{itemize}

\subsection{正交性保证}

\textbf{关键}：约束力与满足约束的速度正交。

\textbf{证明}：

设$V \in \text{Null}(A(\theta))$（满足约束的速度），$F_{\text{constraint}} = A^T(\theta)\lambda$（约束力）。

计算内积：
\begin{align}
F_{\text{constraint}}^T V &= (A^T(\theta)\lambda)^T V \\
&= \lambda^T A(\theta) V \\
&= \lambda^T \cdot 0 = 0
\end{align}

因此，$F_{\text{constraint}}^T V = 0$，即约束力与满足约束的速度正交。$\square$

\textbf{物理意义}：
\begin{itemize}
\item 约束力不会对满足约束的速度做功
\item 约束力只影响违反约束的速度
\item 这确保了约束力不会影响自由运动方向
\end{itemize}

\section{几何直观}

\subsection{6维空间的分解}

\textbf{6维任务空间}可以分解为两个正交的子空间：

\begin{enumerate}
\item \textbf{约束子空间}（维度$6-k$）：
\begin{itemize}
\item 由$\text{Null}(A(\theta))$定义
\item 速度$V$可以自由运动的方向
\item 约束力在这个子空间中为零
\end{itemize}

\item \textbf{约束力子空间}（维度$k$）：
\begin{itemize}
\item 由$\text{Row}(A(\theta))$定义
\item 约束力$F_{\text{constraint}}$必须位于的方向
\item 这是约束子空间的正交补空间
\end{itemize}
\end{enumerate}

\subsection{投影}

\textbf{投影矩阵}：
\begin{equation}
P = I - A^T(A\Lambda^{-1}A^T)^{-1}A\Lambda^{-1}
\end{equation}

\textbf{作用}：
\begin{itemize}
\item $P$将向量投影到约束子空间（$\text{Null}(A)$）
\item $I - P$将向量投影到约束力子空间（$\text{Row}(A)$）
\end{itemize}

\section{总结}

\subsection{零空间}

\begin{enumerate}
\item \textbf{定义}：$\text{Null}(A) = \{x : Ax = 0\}$
\item \textbf{维度}：$n - r$（$r$是$A$的秩）
\item \textbf{在约束问题中}：表示满足约束的速度子空间
\item \textbf{物理意义}：可以自由运动的方向
\end{enumerate}

\subsection{行空间}

\begin{enumerate}
\item \textbf{定义}：$\text{Row}(A) = \text{Col}(A^T)$
\item \textbf{维度}：$r$（$A$的秩）
\item \textbf{在约束问题中}：表示约束力子空间
\item \textbf{物理意义}：约束力必须位于的方向
\end{enumerate}

\subsection{关系}

\begin{enumerate}
\item \textbf{正交性}：$\text{Null}(A) \perp \text{Row}(A)$
\item \textbf{互补性}：$\dim(\text{Null}(A)) + \dim(\text{Row}(A)) = n$
\item \textbf{分解}：任何向量可以唯一分解为零空间分量和行空间分量
\end{enumerate}

\subsection{在约束力中的应用}

\begin{enumerate}
\item \textbf{约束子空间}：$\text{Null}(A(\theta))$（速度必须位于）
\item \textbf{约束力子空间}：$\text{Row}(A(\theta))$（约束力必须位于）
\item \textbf{约束力形式}：$F_{\text{constraint}} = A^T(\theta)\lambda \in \text{Row}(A(\theta))$
\item \textbf{正交性}：约束力与满足约束的速度正交
\end{enumerate}

\end{document}

